% sec1_5.tex — Section 1.5: Relation to Maximum Likelihood

Having specified the distribution of the error vector $\bm{\varepsilon}$, we can use the
\textbf{maximum likelihood (ML) principle} to estimate the model parameters
$(\bm{\beta}, \sigma^2)$.\footnote{For a fuller treatment of maximum likelihood, see Chapter~7.}
In this section, we will show that $\vec{b}$, the OLS estimator of $\bm{\beta}$, is also the
ML estimator, and the OLS estimator of $\sigma^2$ differs only slightly from the ML
counterpart, when the error is normally distributed.  We will also show that $\vec{b}$
achieves the \textbf{Cramer-Rao lower bound}.

%% ─── The Maximum Likelihood Principle ─────────────────────────────────────────
\subsection*{The Maximum Likelihood Principle}

As you might recall from elementary statistics, the basic idea of the ML principle is to
choose the parameter estimates to maximize the probability of obtaining the observed
sample.  To be more precise, we assume that the probability density of the sample
$(\vec{y}, \vec{X})$ is a member of a family of functions indexed by a finite-dimensional
parameter vector $\tilde{\bm{\zeta}}$: $f(\vec{y}, \vec{X}; \tilde{\bm{\zeta}})$.
(This is described as \textbf{parameterizing} the density function.)  This function, viewed
as a function of the hypothetical parameter vector $\tilde{\bm{\zeta}}$, is called the
\textbf{likelihood function}.  At the true parameter vector $\bm{\zeta}$, the density of
$(\vec{y}, \vec{X})$ is $f(\vec{y}, \vec{X}; \bm{\zeta})$.  The ML estimate of the true
parameter vector $\bm{\zeta}$ is the $\tilde{\bm{\zeta}}$ that maximizes the likelihood
function given the data $(\vec{y}, \vec{X})$.

%% ─── Conditional versus Unconditional Likelihood ──────────────────────────────
\subsection*{Conditional versus Unconditional Likelihood}

Since a (joint) density is the product of a marginal density and a conditional density,
the density of $(\vec{y}, \vec{X})$ can be written as
%
\begin{equation}
  f(\vec{y}, \vec{X}; \bm{\zeta})
    = f(\vec{y} \mid \vec{X}; \bm{\theta}) \cdot f(\vec{X}; \bm{\psi}),
  \label{eq:1.5.1}
\end{equation}
%
where $\bm{\theta}$ is the subset of the parameter vector $\bm{\zeta}$ that determines the
conditional density function and $\bm{\psi}$ is the subset determining the marginal density
function.  The parameter vector of interest is $\bm{\theta}$; for the linear regression model
with normal errors, $\bm{\theta} = (\bm{\beta}', \sigma^2)'$ and $f(\vec{y} \mid \vec{X};
\bm{\theta})$ is given by \eqref{eq:1.5.4} below.

Let $\tilde{\bm{\zeta}} \equiv (\tilde{\bm{\theta}}', \tilde{\bm{\psi}}')'$ be a hypothetical
value of $\bm{\zeta} = (\bm{\theta}', \bm{\psi}')'$.  Then the (unconditional or joint)
likelihood function is
%
\begin{equation}
  f(\vec{y}, \vec{X}; \tilde{\bm{\zeta}})
    = f(\vec{y} \mid \vec{X}; \tilde{\bm{\theta}})
      \cdot f(\vec{X}; \tilde{\bm{\psi}}).
  \label{eq:1.5.2}
\end{equation}
%
If we knew the parametric form of $f(\vec{X}; \tilde{\bm{\psi}})$, then we could maximize
this joint likelihood function over the entire hypothetical parameter vector
$\tilde{\bm{\zeta}}$, and the ML estimate of $\bm{\theta}$ would be the elements of the ML
estimate of $\bm{\zeta}$.  We cannot do this for the classical regression model because the
model does not specify $f(\vec{X}; \tilde{\bm{\psi}})$.  However, if there is no functional
relationship between $\tilde{\bm{\theta}}$ and $\tilde{\bm{\psi}}$ (such as a subset of
$\tilde{\bm{\psi}}$ being a function of $\tilde{\bm{\theta}}$), then maximizing
\eqref{eq:1.5.2} with respect to $\tilde{\bm{\zeta}}$ is achieved by separately maximizing
$f(\vec{y} \mid \vec{X}; \tilde{\bm{\theta}})$ with respect to $\tilde{\bm{\theta}}$ and
maximizing $f(\vec{X}; \tilde{\bm{\psi}})$ with respect to $\tilde{\bm{\psi}}$.  Thus the ML
estimate of $\bm{\theta}$ also maximizes the \textbf{conditional likelihood}
$f(\vec{y} \mid \vec{X}; \tilde{\bm{\theta}})$.

%% ─── The Log Likelihood for the Regression Model ─────────────────────────────
\subsection*{The Log Likelihood for the Regression Model}

As already observed, Assumption~1.5 (the normality assumption) together with
Assumptions~1.2 and~1.4 imply that the distribution of $\bm{\varepsilon}$ conditional on
$\vec{X}$ is $N(\vec{0}, \sigma^2 \vec{I}_n)$ (see (1.4.1)).  But since
$\vec{y} = \vec{X}\bm{\beta} + \bm{\varepsilon}$ by Assumption~1.1, we have
%
\begin{equation}
  \vec{y} \mid \vec{X} \sim N(\vec{X}\bm{\beta},\; \sigma^2 \vec{I}_n).
  \label{eq:1.5.3}
\end{equation}
%
Thus, the conditional density of $\vec{y}$ given $\vec{X}$
is\footnote{Recall from basic probability theory that the density function for an
$n$-variate normal distribution with mean $\bm{\mu}$ and variance matrix $\bm{\Sigma}$ is
\[
  (2\pi)^{-n/2} \lvert \bm{\Sigma} \rvert^{-1/2}
  \exp\!\biggl[-\frac{1}{2}(\vec{y} - \bm{\mu})' \bm{\Sigma}^{-1}
  (\vec{y} - \bm{\mu})\biggr].
\]
To derive \eqref{eq:1.5.4}, just set $\bm{\mu} = \vec{X}\bm{\beta}$ and
$\bm{\Sigma} = \sigma^2 \vec{I}_n$.}
%
\begin{equation}
  f(\vec{y} \mid \vec{X})
    = (2\pi\sigma^2)^{-n/2}
      \exp\!\biggl[-\frac{1}{2\sigma^2}
      (\vec{y} - \vec{X}\bm{\beta})'(\vec{y} - \vec{X}\bm{\beta})\biggr].
  \label{eq:1.5.4}
\end{equation}
%
Replacing the true parameters $(\bm{\beta}, \sigma^2)$ by their hypothetical values
$(\widetilde{\bm{\beta}}, \tilde{\sigma}^2)$ and taking logs, we obtain the
\textbf{log likelihood function}:
%
\begin{equation}
  \log L(\widetilde{\bm{\beta}}, \tilde{\sigma}^2)
    = -\frac{n}{2}\log(2\pi)
      - \frac{n}{2}\log(\tilde{\sigma}^2)
      - \frac{1}{2\tilde{\sigma}^2}
        (\vec{y} - \vec{X}\widetilde{\bm{\beta}})'
        (\vec{y} - \vec{X}\widetilde{\bm{\beta}}).
  \label{eq:1.5.5}
\end{equation}
%
Since the log transformation is a monotone transformation, the ML estimator of
$(\bm{\beta}, \sigma^2)$ is the $(\widetilde{\bm{\beta}}, \tilde{\sigma}^2)$ that maximizes
this log likelihood.

%% ─── ML via Concentrated Likelihood ──────────────────────────────────────────
\subsection*{ML via Concentrated Likelihood}

It is instructive to maximize the log likelihood in two stages.  First, maximize over
$\widetilde{\bm{\beta}}$ for any given $\tilde{\sigma}^2$.  The $\widetilde{\bm{\beta}}$
that maximizes the objective function could (but does not, in the present case of
Assumptions~1.1--1.5) depend on $\tilde{\sigma}^2$.  Second, maximize over
$\tilde{\sigma}^2$ taking into account that the $\widetilde{\bm{\beta}}$ obtained in the
first stage could depend on $\tilde{\sigma}^2$.  The log likelihood function in which
$\widetilde{\bm{\beta}}$ is constrained to be the value from the first stage is called the
\textbf{concentrated log likelihood function} (concentrated with respect to
$\widetilde{\bm{\beta}}$).  For the normal log likelihood \eqref{eq:1.5.5}, the first stage
amounts to minimizing the sum of squares
$(\vec{y} - \vec{X}\widetilde{\bm{\beta}})'(\vec{y} - \vec{X}\widetilde{\bm{\beta}})$.
The $\widetilde{\bm{\beta}}$ that does it is none other than the OLS estimator $\vec{b}$,
and the minimized sum of squares is $\vec{e}'\vec{e}$.  Thus, the concentrated log
likelihood is
%
\begin{equation}
  \text{concentrated log likelihood}
    = -\frac{n}{2}\log(2\pi)
      - \frac{n}{2}\log(\tilde{\sigma}^2)
      - \frac{1}{2\tilde{\sigma}^2}\,\vec{e}'\vec{e}.
  \label{eq:1.5.6}
\end{equation}
%
This is a function of $\tilde{\sigma}^2$ alone, and the $\tilde{\sigma}^2$ that maximizes the
concentrated likelihood is the ML estimate of $\sigma^2$.  The maximization is
straightforward for the present case of the classical regression model, because
$\vec{e}'\vec{e}$ is not a function of $\tilde{\sigma}^2$ and so can be taken as a constant.
Still, taking the derivative with respect to $\tilde{\sigma}^2$, rather than with respect to
$\tilde{\sigma}$, can be tricky.  This can be avoided by denoting $\tilde{\sigma}^2$ by
$\tilde{\gamma}$.  Taking the derivative of \eqref{eq:1.5.6} with respect to
$\tilde{\gamma}$ ($\equiv \tilde{\sigma}^2$) and setting it to zero, we obtain the following
result.

\begin{proposition}[ML Estimator of $(\bm{\beta}, \sigma^2)$]
\label{prop:1.5}
\textit{Suppose Assumptions~1.1--1.5 hold.
Then the ML estimator of\/ $\bm{\beta}$ is the OLS estimator\/ $\vec{b}$ and}
%
\begin{equation}
  \textit{ML estimator of\/ $\sigma^2$}
    = \frac{1}{n}\,\vec{e}'\vec{e}
    = \frac{SSR}{n}
    = \frac{n - K}{n}\,s^2.
  \label{eq:1.5.7}
\end{equation}
\end{proposition}

We know from Proposition~1.2 that $s^2$ is unbiased.  Since $s^2$ is multiplied by a
factor $(n - K)/n$ which is different from~1, the ML estimator of $\sigma^2$ is biased,
although the bias becomes arbitrarily small as the sample size $n$ increases for any given
fixed $K$.

For later use, we calculate the maximized value of the likelihood function.  Substituting
\eqref{eq:1.5.7} into \eqref{eq:1.5.6}, we obtain
%
\[
  \text{maximized log likelihood}
    = -\frac{n}{2}\log\!\biggl(\frac{2\pi}{n}\biggr)
      - \frac{n}{2}
      - \frac{n}{2}\log(SSR),
\]
%
so that the maximized likelihood is
%
\begin{equation}
  \max_{\widetilde{\bm{\beta}},\, \tilde{\sigma}^2}
    L(\widetilde{\bm{\beta}}, \tilde{\sigma}^2)
    = \biggl(\frac{2\pi}{n}\biggr)^{-n/2}
      \cdot \exp\!\biggl(-\frac{n}{2}\biggr)
      \cdot (SSR)^{-n/2}.
  \label{eq:1.5.8}
\end{equation}

%% ─── Cramer-Rao Bound for the Classical Regression Model ─────────────────────
\subsection*{Cramer-Rao Bound for the Classical Regression Model}

Just to refresh your memory of basic statistics, we temporarily step outside the classical
regression model and present without proof the Cramer-Rao inequality for the
variance-covariance matrix of any unbiased estimator.  For this purpose, define the
\textbf{score vector} at a hypothetical parameter value $\tilde{\bm{\theta}}$ to be the
gradient (vector of partial derivatives) of log likelihood:
%
\begin{equation}
  \text{score:}\quad
  \vec{s}(\tilde{\bm{\theta}})
    \equiv \frac{\partial \log L(\tilde{\bm{\theta}})}
                {\partial \tilde{\bm{\theta}}}.
  \label{eq:1.5.9}
\end{equation}

\medskip
\noindent\textbf{Cramer-Rao Inequality:}\enspace
\textit{Let\/ $\vec{z}$ be a vector of random variables (not necessarily independent) the
joint density of which is given by $f(\vec{z}; \bm{\theta})$, where $\bm{\theta}$ is an
$m$-dimensional vector of parameters in some parameter space $\Theta$.  Let
$L(\tilde{\bm{\theta}}) \equiv f(\vec{z}; \tilde{\bm{\theta}})$ be the likelihood function,
and let\/ $\hat{\bm{\theta}}(\vec{z})$ be an unbiased estimator of\/ $\bm{\theta}$ with a
finite variance-covariance matrix.  Then, under some regularity conditions on
$f(\vec{z}; \bm{\theta})$ (not stated here),}
%
\[
  \Var[\hat{\bm{\theta}}(\vec{z})]
    \underset{(m \times m)}{\geq}
    \vec{I}(\bm{\theta})^{-1}
  \quad (\equiv \textit{Cramer-Rao Bound}),
\]
%
\textit{where\/ $\vec{I}(\bm{\theta})$ is the \textbf{information matrix} defined by}
%
\begin{equation}
  \vec{I}(\bm{\theta})
    \equiv \E[\vec{s}(\bm{\theta})\,\vec{s}(\bm{\theta})'].
  \label{eq:1.5.10}
\end{equation}
%
\textit{(Note well that the score is evaluated at the true parameter value $\bm{\theta}$.)
Also under the regularity conditions, the information matrix equals the negative of the
expected value of the Hessian (matrix of second partial derivatives) of the log likelihood:}
%
\begin{equation}
  \vec{I}(\bm{\theta})
    = -\,\E\!\biggl[
        \frac{\partial^2 \log L(\bm{\theta})}
             {\partial \tilde{\bm{\theta}}\,\partial \tilde{\bm{\theta}}'}
      \biggr].
  \label{eq:1.5.11}
\end{equation}
%
\textit{This is called the \textbf{information matrix equality}.}

\medskip

See, e.g., Amemiya (1985, Theorem~1.3.1) for a proof and a statement of the regularity
conditions.  Those conditions guarantee that the operations of differentiation and taking
expectations can be interchanged.  Thus, for example,
%
\[
  \E[\partial L(\bm{\theta}) / \partial \tilde{\bm{\theta}}]
    = \partial\, \E[L(\bm{\theta})] / \partial \tilde{\bm{\theta}}.
\]

Now, for the classical regression model (of Assumptions~1.1--1.5), the likelihood function
$L(\tilde{\bm{\theta}})$ in the Cramer-Rao inequality is the conditional density
\eqref{eq:1.5.4}, so the variance in the inequality is the variance conditional on
$\vec{X}$.  It can be shown that those regularity conditions are satisfied for the normal
density \eqref{eq:1.5.4} (see, e.g., Amemiya, 1985, Sections~1.3.2 and~1.3.3).  In the
rest of this subsection, we calculate the information matrix for \eqref{eq:1.5.4}.  The
parameter vector $\bm{\theta}$ is $(\bm{\beta}', \sigma^2)'$.

So $\tilde{\bm{\theta}} = (\widetilde{\bm{\beta}}', \tilde{\gamma})'$ and the matrix of
second derivatives we seek to calculate is
%
\begin{equation}
  \frac{\partial^2 \log L(\bm{\theta})}
       {\partial \tilde{\bm{\theta}}\,\partial \tilde{\bm{\theta}}'}
  \underset{((K+1) \times (K+1))}{=}
  \begin{bmatrix}
    \dfrac{\partial^2 \log L(\bm{\theta})}
          {\partial \widetilde{\bm{\beta}}\,\partial \widetilde{\bm{\beta}}'}
    & \dfrac{\partial^2 \log L(\bm{\theta})}
            {\partial \widetilde{\bm{\beta}}\,\partial \tilde{\gamma}} \\[12pt]
    \underset{(K \times K)}{}
    & \underset{(K \times 1)}{} \\[4pt]
    \dfrac{\partial^2 \log L(\bm{\theta})}
          {\partial \tilde{\gamma}\,\partial \widetilde{\bm{\beta}}'}
    & \dfrac{\partial^2 \log L(\bm{\theta})}
            {\partial \tilde{\gamma}^2} \\[6pt]
    \underset{(1 \times K)}{}
    & \underset{(1 \times 1)}{}
  \end{bmatrix}.
  \label{eq:1.5.12}
\end{equation}

The first and second derivatives of the log likelihood \eqref{eq:1.5.5} with respect to
$\tilde{\bm{\theta}}$, evaluated at the true parameter vector $\bm{\theta}$, are
%
\begin{align}
  \frac{\partial \log L(\bm{\theta})}{\partial \widetilde{\bm{\beta}}}
    &= \frac{1}{\gamma}\,\vec{X}'(\vec{y} - \vec{X}\bm{\beta}),
  \label{eq:1.5.13a} \\[6pt]
  \frac{\partial \log L(\bm{\theta})}{\partial \tilde{\gamma}}
    &= -\frac{n}{2\gamma}
       + \frac{1}{2\gamma^2}
         (\vec{y} - \vec{X}\bm{\beta})'(\vec{y} - \vec{X}\bm{\beta}).
  \label{eq:1.5.13b}
\end{align}
%
\begin{align}
  \frac{\partial^2 \log L(\bm{\theta})}
       {\partial \widetilde{\bm{\beta}}\,\partial \widetilde{\bm{\beta}}'}
    &= -\frac{1}{\gamma}\,\vec{X}'\vec{X},
  \label{eq:1.5.14a} \\[6pt]
  \frac{\partial^2 \log L(\bm{\theta})}{\partial \tilde{\gamma}^2}
    &= \frac{n}{2\gamma^2}
       - \frac{1}{\gamma^3}
         (\vec{y} - \vec{X}\bm{\beta})'(\vec{y} - \vec{X}\bm{\beta}),
  \label{eq:1.5.14b} \\[6pt]
  \frac{\partial^2 \log L(\bm{\theta})}
       {\partial \widetilde{\bm{\beta}}\,\partial \tilde{\gamma}}
    &= -\frac{1}{\gamma^2}\,\vec{X}'(\vec{y} - \vec{X}\bm{\beta}).
  \label{eq:1.5.14c}
\end{align}

Since the derivatives are evaluated at the true parameter value,
$\vec{y} - \vec{X}\bm{\beta} = \bm{\varepsilon}$ in these expressions.  Substituting
\eqref{eq:1.5.14a}--\eqref{eq:1.5.14c} into \eqref{eq:1.5.12} and using
$\E(\bm{\varepsilon} \mid \vec{X}) = \vec{0}$ (Assumption~1.2),
$\E(\bm{\varepsilon}'\bm{\varepsilon} \mid \vec{X}) = n\sigma^2$ (implication of
Assumption~1.4), and recalling $\gamma = \sigma^2$, we can easily derive
%
\begin{equation}
  \vec{I}(\bm{\theta})
    = \begin{bmatrix}
        \dfrac{1}{\sigma^2}\,\vec{X}'\vec{X} & \vec{0} \\[8pt]
        \vec{0}' & \dfrac{n}{2\sigma^4}
      \end{bmatrix}.
  \label{eq:1.5.15}
\end{equation}
%
Here, the expectation is conditional on $\vec{X}$ because the likelihood function
\eqref{eq:1.5.4} is a conditional density conditional on $\vec{X}$.  This block diagonal
matrix can be inverted to obtain the Cramer-Rao bound:
%
\begin{equation}
  \text{Cramer-Rao bound}
    \equiv \vec{I}(\bm{\theta})^{-1}
    = \begin{bmatrix}
        \sigma^2 \cdot (\vec{X}'\vec{X})^{-1} & \vec{0} \\[8pt]
        \vec{0}' & \dfrac{2\sigma^4}{n}
      \end{bmatrix}.
  \label{eq:1.5.16}
\end{equation}
%
Therefore, the unbiased estimator $\vec{b}$, whose variance is
$\sigma^2 \cdot (\vec{X}'\vec{X})^{-1}$ by Proposition~1.1, attains the Cramer-Rao bound.
We have thus proved

\begin{proposition}[$\vec{b}$ is the Best Unbiased Estimator (BUE)]
\label{prop:1.6}
\textit{Under Assumptions~1.1--1.5, the OLS estimator\/ $\vec{b}$ of\/ $\bm{\beta}$ is BUE
in that any other unbiased (but not necessarily linear) estimator has larger conditional
variance in the matrix sense.}
\end{proposition}

This result should be distinguished from the Gauss-Markov Theorem that $\vec{b}$ is
minimum variance among those estimators that are unbiased \textit{and} linear in $\vec{y}$.
Proposition~1.6 says that $\vec{b}$ is minimum variance in a larger class of estimators
that includes nonlinear unbiased estimators.  This stronger statement is obtained under the
normality assumption (Assumption~1.5) which is not assumed in the Gauss-Markov Theorem.
Put differently, the Gauss-Markov Theorem does not exclude the possibility of some
nonlinear estimator beating OLS, but this possibility is ruled out by the normality
assumption.

As was already seen, the ML estimator of $\sigma^2$ is biased, so the Cramer-Rao bound
does not apply.  But the OLS estimator $s^2$ of $\sigma^2$ is unbiased.  Does it achieve
the bound?  We have shown in a review question to the previous section that
%
\[
  \Var(s^2 \mid \vec{X}) = \frac{2\sigma^4}{n - K}
\]
%
under the same set of assumptions as in Proposition~1.6.  Therefore, $s^2$ does not attain
the Cramer-Rao bound $2\sigma^4/n$.  However, it can be shown that an unbiased estimator of
$\sigma^2$ with variance lower than $2\sigma^4/(n - K)$ does not exist (see, e.g., Rao,
1973, p.~319).

%% ─── The F-Test as a Likelihood Ratio Test ───────────────────────────────────
\subsection*{The $F$-Test as a Likelihood Ratio Test}

The \textbf{likelihood ratio test} of the null hypothesis compares $L_U$, the maximized
likelihood without the imposition of the restriction specified in the null hypothesis, with
$L_R$, the likelihood maximized subject to the restriction.  If the likelihood ratio
$\lambda \equiv L_U / L_R$ is too large, it should be a sign that the null is false.  The
$F$-test of the null hypothesis $\mathrm{H}_0\colon \vec{R}\bm{\beta} = \vec{r}$ considered
in the previous section is a likelihood ratio test because the $F$-ratio is a monotone
transformation of the likelihood ratio $\lambda$.  For the present model, $L_U$ is given by
\eqref{eq:1.5.8} where the $SSR$, the sum of squared residuals minimized without the
constraint $\mathrm{H}_0$, is the $SSR_U$ in (1.4.11).  The restricted likelihood $L_R$ is
given by replacing this $SSR$ by the restricted sum of squared residuals, $SSR_R$.  So
%
\begin{equation}
  L_R
    = \max_{\widetilde{\bm{\beta}},\,\tilde{\sigma}^2
            \;\text{s.t.}\;\mathrm{H}_0}
      L(\widetilde{\bm{\beta}}, \tilde{\sigma}^2)
    = \biggl(\frac{2\pi}{n}\biggr)^{-n/2}
      \cdot \exp\!\biggl(-\frac{n}{2}\biggr)
      \cdot (SSR_R)^{-n/2},
  \label{eq:1.5.17}
\end{equation}
%
and the likelihood ratio is
%
\begin{equation}
  \lambda
    \equiv \frac{L_U}{L_R}
    = \biggl(\frac{SSR_U}{SSR_R}\biggr)^{-n/2}.
  \label{eq:1.5.18}
\end{equation}
%
Comparing this with the formula (1.4.11) for the $F$-ratio, we see that the $F$-ratio is a
monotone transformation of the likelihood ratio $\lambda$:
%
\begin{equation}
  F = \frac{n - K}{\#\vec{r}}\,(\lambda^{2/n} - 1),
  \label{eq:1.5.19}
\end{equation}
%
so that the two tests are the same.

%% ─── Quasi-Maximum Likelihood ────────────────────────────────────────────────
\subsection*{Quasi-Maximum Likelihood}

All these results assume the normality of the error term.  Without normality, there is no
guarantee that the ML estimator of $\bm{\beta}$ is OLS (Proposition~1.5) or that the OLS
estimator $\vec{b}$ achieves the Cramer-Rao bound (Proposition~1.6).  However,
Proposition~1.5 does imply that $\vec{b}$ is a \textbf{quasi-} (or \textbf{pseudo-)
maximum likelihood estimator}, an estimator that maximizes a misspecified likelihood
function.  The misspecified likelihood function we have considered is the normal likelihood.
The results of Section~1.3 can then be interpreted as providing the finite-sample properties
of the quasi-ML estimator when the error is incorrectly specified to be normal.
